\section*{Capítulo \ref{ch:g10:chemical-species} --- Espécies químicas naturais e sintéticas}

\begin{solution}{exo:g10:chemical-species:1}
Naturais: a vanilina de uma vagem, a cafeína dos grãos de café.
Sintéticas: a aspirina, a vanilina de fábrica.
\end{solution}

\begin{solution}{exo:g10:chemical-species:2}
Elas são a mesma \omterm{def:g6:pure-substances-and-mixtures:species}{espécie química}, e uma espécie tem as mesmas
propriedades seja qual for a sua origem.
\end{solution}

\begin{solution}{exo:g10:chemical-species:3}
$R_f = 2.4/6.0 = 0.40$.
\end{solution}

\begin{solution}{exo:g10:chemical-species:4}
Ele resfria o vapor e o faz voltar a ser líquido. Entrando pela
extremidade de baixo, a água enche toda a camisa (e encontra o vapor mais
quente por último, resfriando melhor).
\end{solution}

\begin{solution}{exo:g10:chemical-species:5}
\omterm{def:g2:mixing-and-dissolving:dissolve}{Dissolver} a espécie muito melhor do que a água; ser \omterm{def:g6:solutions-and-solubility:miscible}{imiscível} com a
água; ser o menos perigoso possível e fácil de retirar.
\end{solution}

\begin{solution}{exo:g10:chemical-species:6}
O etanol é \omterm{def:g6:solutions-and-solubility:miscible}{miscível} com a água: não se formam camadas, e nada pode ser
separado. O ciclo-hexano é \omterm{def:g6:solutions-and-solubility:miscible}{imiscível} com a água e dissolve bem o linalol.
\end{solution}

\begin{solution}{exo:g10:chemical-species:7}
A água fica em cima: \qty{1}{mL} de diclorometano pesa \qty{1.33}{g},
mais do que \qty{1}{mL} de água, e por isso ele afunda.
\end{solution}

\begin{solution}{exo:g10:chemical-species:8}
Ele funde baixo demais (aspirina: \qty{135}{\celsius}) e ao longo de uma
faixa: a amostra é impura, ou não é aspirina.
\end{solution}

\begin{solution}{exo:g10:chemical-species:9}
Linalol e etanoato de linalila (as duas manchas dele estão nas alturas
das duas referências). O etanoato de linalila subiu mais alto.
\end{solution}

\begin{solution}{exo:g10:chemical-species:10}
\omterm{def:g10:chemical-species:distillation}{Hidrodestilação}; acréscimo de ciclo-hexano ao destilado e agitação;
separação das camadas; \omterm{def:g3:separating-mixtures:evaporate}{evaporação} do ciclo-hexano.
\end{solution}

\begin{solution}{exo:g10:chemical-species:11}
O grafite não \omterm{def:g2:mixing-and-dissolving:dissolve}{se dissolve} no \omterm{def:g10:chemical-species:tlc}{eluente}; acima do \omterm{def:g10:chemical-species:tlc}{eluente}, as manchas são
levadas para cima pela placa em vez de se dissolverem na cuba.
\end{solution}

\begin{solution}{exo:g10:chemical-species:12}
O ciclo-hexano ferve a \qty{81}{\celsius}, o linalol a
\qty{198}{\celsius}: um aquecimento suave expulsa o ciclo-hexano
enquanto o linalol fica.
\end{solution}

\begin{solution}{exo:g10:chemical-species:13}
A amostra provavelmente é essa espécie (mesmo $R_f$ que a natural na
mesma placa) e provavelmente é pura (uma única mancha, ponto de fusão
nítido).
\end{solution}

\begin{solution}{exo:g10:chemical-species:14}
$0.40 \times 5.5 = \qty{2.2}{cm}$. A outra está a $0.75 \times 5.5 =
\qty{4.1}{cm}$ (arredondando ao milímetro), portanto cerca de
\qty{1.9}{cm} acima de A.
\end{solution}

\begin{solution}{exo:g10:chemical-species:15}
\qty{200}{mL} podem \omterm{def:g2:mixing-and-dissolving:dissolve}{dissolver} $1.6 \times 0.200 = \qty{0.32}{g}$: com
\qty{0.20}{g}, ele não está saturado; ainda poderia \omterm{def:g2:mixing-and-dissolving:dissolve}{dissolver}
\qty{0.12}{g}. Os \qty{0.20}{g} dissolvidos seriam perdidos com a água;
o ciclo-hexano, que dissolve muito melhor o linalol, recupera a maior
parte deles.
\end{solution}

\begin{solution}{pb:g10:chemical-species:1}
\textbf{1.} O vapor dela arrasta as espécies perfumadas para fora das
flores e até o condensador.

\textbf{2.} Ele resfria o vapor e o faz voltar a ser líquido; a água fria
entra pela extremidade de baixo.

\textbf{3.} \omterm{def:g6:solutions-and-solubility:miscible}{Imiscível}; menos denso que a água (ele flutua).

\textbf{4.} A camada é muito fina, e uma parte do óleo está dissolvida
na água ou dispersa nela: uma \omterm{def:g10:chemical-species:extraction}{extração} a recupera.

\textbf{5.} Ele dissolve bem o linalol, é \omterm{def:g6:solutions-and-solubility:miscible}{imiscível} com a água e ferve a
uma temperatura baixa (\qty{81}{\celsius}), por isso é fácil de retirar.

\textbf{6.} Não: o etanol se mistura com a água e não se formam camadas.

\textbf{7.} Em cima (\qty{0.78}{g} por mL, menos que a água). Com o
diclorometano (\qty{1.33}{g} por mL), a camada de \omterm{def:g6:solutions-and-solubility:solute}{solvente} ficaria
embaixo.

\textbf{8.} Nenhuma chama por perto (inflamável); trabalhar numa capela
de exaustão, com luvas, e recolher os resíduos em vez de jogá-los na pia
(tóxico para os organismos aquáticos).

\textbf{9.} O ciclo-hexano ferve a \qty{81}{\celsius}, o linalol a
\qty{198}{\celsius}: um aquecimento suave só retira o \omterm{def:g6:solutions-and-solubility:solute}{solvente}.

\textbf{10.} $R_f(\text{L}) = 2.4/6.0 = 0.40$; $R_f(\text{A}) = 4.5/6.0
= 0.75$.

\textbf{11.} Linalol e etanoato de linalila.

\textbf{12.} Não: ele contém pelo menos duas espécies; é uma \omterm{def:g6:pure-substances-and-mixtures:mixture}{mistura}.

\textbf{13.} Para que todas as manchas sejam eluídas nas mesmas
condições: só assim as alturas podem ser comparadas.

\textbf{14.} O $R_f$ dele é igual ao do linalol nas mesmas condições: ele
provavelmente é linalol, idêntico à \omterm{def:g10:chemical-species:natural}{espécie natural}.

\textbf{15.} $1.0 \times 0.90 = \qty{0.90}{g}$.

\textbf{16.} $0.90/100 = \qty{0.90}{\percent}$.

\textbf{17.} $100/0.009 \approx 11\,000$\,g, cerca de \qty{11}{kg} de
flores.

\textbf{18.} Ele é extraído de uma planta, sem \omterm{def:g5:irreversible-changes:chemical-change}{transformação química}. Mas
o linalol dele é a mesma espécie que o linalol sintético: mesma fórmula,
mesmas propriedades.

\textbf{19.} \qty{0.90}{g} de óleo essencial a partir de \qty{100}{g} de
flores.
\end{solution}
